\chapter{Elliptic Thom lines and derived incidence}
\label{ei:chap:incidence}

The elliptic class of Chapter~\ref{ed:chap:euler} transports
cohomological correspondences. A sheaf-valued coefficient requires
the line carrying that class and the map defined by its section.
Already for two interacting colors, the line has nonzero degree.
Its section vanishes twice at collision. A change of line objects
can remove a tensor twist while preserving this doubled zero.

We follow the pair coefficient into its simultaneous derived
intersections. A spanning tree reduces each connected graph to a
finite algebra. The first cycle gives a nonzero ternary operation,
and a repeated color changes the character of its perfect trace.

We then study two parameter families of these algebras. Independent
quadratic parameters produce discriminant modules. The actual
elliptic displacement gives a different family, with exceptional
fibers over the elliptic curve and a square root on a halving cover.
These calculations specify coefficient data for the compact
comparison posed at the end of the chapter.

\section{The pair coefficient and its image}
\label{ei:sec:pair}

Fix a smooth complex elliptic curve $E$ with origin, and put
$\Theta=\mathcal O_E(0)$ with divisor section $\vartheta$.
This $E$ parametrizes spectral variables. Its role differs from
the elliptic factor of a threefold and from a curve of field insertions.

Take one point of each of two distinct colors $i,j$, with two
arrows from $i$ to $j$. Choose weights $(2,0)$ and $(0,2)$.
Their compatibility imposes
$2t_1=2t_2$ on the elliptic equivariant parameters.
In the formal neighborhood of the origin this forces $t_1=t_2$.

We work on the component $t_1=t_2=t$ and set $h=2t$.
With $u=z_j-z_i$, the forward arrow factor is
\begin{equation}\label{ei:eq:pair}
 s_h(u)=\vartheta(u)\vartheta(u+h),\qquad
 \mathcal S=u^*\Theta\otimes(u+h)^*\Theta,\qquad
 \mathcal S^{-1}\xrightarrow{s_h}\mathcal O.
\end{equation}
The maps $u$ and $u+h$ have domain $E_u\times E_h$.
There is no same-color factor for this dimension vector.

These are the compatible weights and arrow factor of
Yang--Zhao.\footnote{Y.~Yang and G.~Zhao,
\emph{Quiver varieties and elliptic quantum groups},
\href{https://arxiv.org/abs/1708.01418}{Assumption~2.1,
Remark~2.2(2) and equations~(2)--(3)}.
Their Theorem~8.4 compares algebras of rational sections.
The isomorphism in Remark~8.5 requires localization.}

At $h=0$, write $d:E^2\to E$ for the difference map and
$\Delta=d^{-1}(0)$. The coefficient map becomes
\begin{equation}\label{ei:eq:doubled-pair}
 \mathcal L=d^*\Theta^{-2}
    =\mathcal O_{E^2}(-2\Delta)
 \xrightarrow{\vartheta(d)^2}\mathcal O_{E^2}.
\end{equation}
Its image is the ideal of $2\Delta$, and its cokernel is
$\mathcal O_{2\Delta}$.

\begin{proposition}\label{ei:prop:jet}
The line $\mathcal L$ is nontrivial. On a normal formal disk to
$\Delta$, its coefficient cokernel is $\mathbb C[[u]]/(u^2)$.
The sequence
\begin{equation}\label{ei:eq:jet}
 0\longrightarrow\mathcal O_\Delta(-\Delta)
 \longrightarrow\mathcal O_{2\Delta}
 \longrightarrow\mathcal O_\Delta\longrightarrow0
\end{equation}
does not split over the ambient local ring.
\end{proposition}

\begin{proof}
Restrict $\mathcal L$ to $\{z_1\}\times E$. Its degree is $-2$,
whereas a trivial line has degree zero.
Choose an invariant differential on $E$ and its normalized formal
logarithm $u$. In a theta frame,
$\vartheta(u)=u\,v(u)$ with $v(0)\ne0$.
Rescaling the source frame removes $v(u)^2$, leaving the map $u^2$.

The ideal $(u)/(u^2)$ is the first term of \eqref{ei:eq:jet}.
A module section would lift $1$ to an element with nonzero constant
term annihilated by $u$. Such an element does not exist in
$\mathbb C[[u]]/(u^2)$.
\end{proof}

Thus restriction to reduced support loses a normal extension class.
Even a regular scalar change of frame preserves its doubled vanishing.
At nonzero displacement, the two orders also have different images.
In an odd theta frame the reversed factor is
$\vartheta(u)\vartheta(u-h)$.
Writing $D_a=\{u=a\}$ on $E_u$, the two image ideals are
\[
 \mathcal I_+=\mathcal O(-D_0-D_{-h}),\qquad
 \mathcal I_-=\mathcal O(-D_0-D_h).
\]
Their sum is the image of the two maps together.

\begin{proposition}\label{ei:prop:ordered-images}
On a fixed fiber, $\mathcal J_h=\mathcal I_++\mathcal I_-$ is
\[
 \mathcal J_h=
 \begin{cases}
  \mathcal O(-D_0),&2h\ne0,\\
  \mathcal O(-D_0-D_h),&2h=0,\ h\ne0,\\
  \mathcal O(-2D_0),&h=0.
 \end{cases}
\]
Near $(u,h)=(0,0)$ the image ideal is
\begin{equation}\label{ei:eq:sum-ideal}
 J=(u(u+h),u(u-h))=(u^2,uh)\subset\mathbb C[[u,h]].
\end{equation}
Its quotient contains a nonzero class $[u]$ killed by $h$.
Formation of the image fails to commute with specialization at $h=0$.
\end{proposition}

\begin{proof}
The sum of two divisor ideals on a smooth curve has the smaller
vanishing order at every point. The common zeros are precisely those
listed in the three cases, with their displayed multiplicities.

Formal logarithms and invertible theta frames give $u(u+h)$ and
$u(u-h)$. Their sum and difference generate $u^2,uh$ because
$2$ is invertible. Hence $u\notin J$ and $hu\in J$.

Moreover $uh$ has nonzero class in $J/hJ$. An equality $uh=hj$
with $j\in J$ would imply $u=j$, contrary to $u\notin J$.
This class maps to zero in $\mathbb C[[u,h]]/(h)$.
The specialized maps themselves have image $(u^2)$.
\end{proof}

When $2h\ne0$, the quotient $\mathcal J_h/\mathcal I_+$ is supported
at $-h$, and $\mathcal J_h/\mathcal I_-$ is supported at $h$.
At $2h=0$ the two image divisors agree.
This identifies divisor lines. An identification of ordered Hall
operations must also retain their scalar and equivariant data.

\section{Flags, line twists and simultaneous equations}
\label{ei:sec:flags}

Set $h=0$. Let $G$ be a finite simple graph with ordered vertex set $V$.
An edge records an interacting pair of multiplicity two.
For disjoint subsets $A,B\subset V$, define
\[
 D_A=\sum_{\substack{i<j,\ i,j\in A\\\{i,j\}\in E(G)}}\Delta_{ij},
 \qquad
 D_{A,B}=\sum_{\substack{i\in A,\ j\in B\\\{i,j\}\in E(G)}}\Delta_{ij}.
\]
Here $\Delta_{ij}$ is the equality divisor on the product of
spectral curves. Pull all lines to the common product and put
\[
 \mathcal I_A=\mathcal O(-2D_A),\qquad
 \mathcal L_{A,B}=\mathcal O(-2D_{A,B}).
\]

The divisor sections define their embeddings into $\mathcal O$.
Write $\operatorname{Perf}(X)$ for the category of complexes locally
quasi-isomorphic to bounded complexes of finite-rank vector bundles.
For $F\in\operatorname{Perf}(E^A)$ and
$G'\in\operatorname{Perf}(E^B)$, define
$F\star G'=(F\boxtimes G')\otimes\mathcal L_{A,B}$.
Here $\boxtimes$ is the derived external tensor product.
No shuffle symmetrization is included in this definition.

\begin{proposition}\label{ei:prop:flags}
Multiplication of the pair sections factors through
\[
 \mathcal I_A\otimes\mathcal I_B\otimes\mathcal L_{A,B}
       \xrightarrow{\sim}\mathcal I_{A\cup B}.
\]
These isomorphisms are compatible with every ordered flag.
The line identifications
\begin{equation}\label{ei:eq:associator}
 \mathcal L_{A,B}\otimes\mathcal L_{A\cup B,C}
 \cong
 \mathcal L_{B,C}\otimes\mathcal L_{A,B\cup C}
\end{equation}
preserve the sections and all regroupings. In particular,
\[
 \mathcal L_{A,B}
 \cong\mathcal I_{A\cup B}\otimes\mathcal I_A^{-1}
                         \otimes\mathcal I_B^{-1}.
\]
The functors $\mathsf T_A(F)=F\otimes\mathcal I_A^{-1}$
untwist $\star$. Halving the divisors gives compatible square
roots of the line identifications.
\end{proposition}

\begin{proof}
The identity
$D_{A\cup B}=D_A+D_B+D_{A,B}$ separates internal and crossing edges.
It gives the first isomorphism, its section and the coboundary formula.
For three blocks, each side of \eqref{ei:eq:associator} contains
every edge between distinct blocks once.
The same count applies to any finite ordered flag.
Thus every composite is the identification of the same divisor line.

Cancelling the internal lines gives
$\mathsf T_{A\cup B}(F\star G')
\cong\mathsf T_A(F)\boxtimes\mathsf T_B(G')$.
The empty sector has $\mathcal I_\varnothing=\mathcal O$.
Replacing each coefficient $-2$ by $-1$ preserves every divisor
addition identity and gives the square roots.
\end{proof}

The doubled-divisor cokernel belongs to the final embedding
$\mathcal I_{A\cup B}\hookrightarrow\mathcal O$.
Under $\mathsf T$, the two-color coefficient map becomes
$\mathcal O\to\mathcal O(2\Delta)$, with its doubled zero unchanged.
The square roots above concern the specified labeled divisor lines
at $h=0$. Compact orientation lines, same-color factors and
shuffle maps require their own comparisons.

To retain simultaneous equations, fix a commutative ring $R$ and
elements $f_1,\ldots,f_r$. Their Koszul algebra is
\begin{equation}\label{ei:eq:koszul}
 K_R(f_1,\ldots,f_r)
  =\bigl(R\otimes\Lambda(\epsilon_1,\ldots,\epsilon_r),d\bigr),
 \qquad |\epsilon_i|=-1,\quad d\epsilon_i=f_i.
\end{equation}
The differential is a graded derivation of degree one.
All tensor signs below use this cochain grading.

Over $R=\mathbb C[[u]]$, $K_R(u^2)$ has only
$H^0=R/(u^2)$. In $K_R(u,u)$, subtract the first odd generator
from the second. Resolving the remaining regular equation gives
\[
 K_R(u,u)\simeq R/(u)\otimes\Lambda(\epsilon),\qquad d\epsilon=0.
\]
The first construction retains a nilpotent coordinate.
The second retains a negative excess class.

For a triangle put $x=z_2-z_1$, $y=z_3-z_2$ in logarithmic
coordinates. The product coefficient has cone
$K_{\mathbb C[[x,y]]}(x^2y^2(x+y)^2)$, supported on three doubled
lines and concentrated in cohomological degree zero.
The simultaneous pair intersection is
\begin{equation}\label{ei:eq:triangle-koszul}
 K_{\mathbb C[[x,y]]}(x^2,y^2,(x+y)^2).
\end{equation}
It is supported at the origin and has negative cohomology.
This additional incidence object does not measure a failure of
associativity of the original product.

For nonzero $f,g$ in a domain there is instead an exact sequence
\[
 0\longrightarrow R/(f)\xrightarrow{g}R/(fg)
   \longrightarrow R/(g)\longrightarrow0.
\]
Its global version carries the corresponding line twists.
This filtration of a product cone is different from the derived
tensor product of the two cones.

\section{A global spectral coefficient and its finite model}
\label{ei:sec:graph}

Suppose that $G$ is connected, with $n\geq1$ vertices and $m$ edges.
Common translation acts freely on $E^n$, and differences from the
first point identify
\[
 B_n=E^n/E_{\mathrm{diag}}\cong E^{n-1}.
\]
Each equality divisor descends to a divisor $D_e$ on $B_n$.
Write $L_e=\mathcal O_{B_n}(2D_e)$ and let $s_e$ be its divisor
section. Define the commutative differential graded sheaf
\begin{equation}\label{ei:eq:global-koszul}
 \mathscr K_G=\bigotimes_{e\in E(G)}
       [L_e^{-1}\xrightarrow{s_e}\mathcal O_{B_n}],
\end{equation}
with $L_e^{-1}$ in degree $-1$.
Its coefficient algebra is $R\Gamma(B_n,\mathscr K_G)$.

\begin{lemma}\label{ei:lem:completion}
The complex $\mathscr K_G$ is perfect and its cohomology sheaves
have finite length at the origin of $B_n$.
Restriction to the formal completion there computes its derived
global sections as a differential graded algebra.
\end{lemma}

\begin{proof}
Every factor is a two-term complex of line bundles.
Away from their common zero locus, an invertible section makes one
factor contractible. Connectedness identifies the common zero locus
with the origin in $B_n$.
The coherent cohomology sheaves therefore have finite length.

They have no higher cohomology and are unchanged by completion.
The hypercohomology spectral sequence proves that restriction
induces a cohomology isomorphism. Restriction also preserves
multiplication and the differential.
\end{proof}

For $n=1$, this algebra is $\mathbb C$ on a point.
The empty configuration has translation quotient
$[\mathrm{pt}/E]=BE$ and is outside this $n\geq1$ construction.
The spectral quotient also differs from the common translation
torsor on compact sheaf moduli.

Formula \eqref{ei:eq:global-koszul} also defines complexes for
disconnected subgraphs on the same vertex set.
Adding an edge $e$ gives the exact triangle
\begin{equation}\label{ei:eq:edge-triangle}
 \mathscr K_G\otimes L_e^{-1}\xrightarrow{s_e}\mathscr K_G
 \longrightarrow\mathscr K_{G+e}
 \longrightarrow\mathscr K_G\otimes L_e^{-1}[1].
\end{equation}
The middle map is inclusion of the even summand and preserves the
differential graded algebra structure. Locally place the new
generator first and write $a+\epsilon_e b$.
The last map sends this element to $b$ in the shifted term.
Its sign follows from
$d(\epsilon_e b)=s_e b-\epsilon_e db$.

Permuting edges permutes odd generators with exterior signs.
Thus the edge additions and all subset cubes commute on the full
Koszul object, including disconnected intermediate subgraphs.

Choose an oriented spanning tree $T$ and let
$u_1,\ldots,u_{n-1}$ be its logarithmic edge differences.
These are coordinates in the completed local ring of $B_n$.
For a chord $c\notin T$, its difference is
\[
 \ell_c=\sum_a\sigma_{ca}u_a,\qquad
 \sigma_{ca}\in\{0,1,-1\}.
\]
The signs record traversal of the tree path.
Put $\beta=m-n+1$ and
$A_T=\mathbb C[u_1,\ldots,u_{n-1}]/(u_1^2,\ldots,u_{n-1}^2)$.

\begin{theorem}\label{ei:thm:finite-model}
After these coordinate and frame choices,
$R\Gamma(B_n,\mathscr K_G)$ is represented by
\begin{equation}\label{ei:eq:finite-model}
 C_{G,T}=A_T\otimes\Lambda(\epsilon_c:c\notin T),\qquad
 d\epsilon_c=2\sum_{a<b}\sigma_{ca}\sigma_{cb}u_au_b.
\end{equation}
Its dimension is $2^m$ and its degrees lie in $[-\beta,0]$.
Different trees give quasi-isomorphic algebras through the full
edge Koszul complex.
\end{theorem}

\begin{proof}
In a theta frame, each section is the square of its logarithmic
difference times a unit. Divide its Koszul generator by that unit.
The tree equations form the regular sequence
$u_1^2,\ldots,u_{n-1}^2$.
Their Koszul resolution maps multiplicatively to $A_T$.
Tensoring with the finite free chord factors preserves the
quasi-isomorphism.

In $A_T$ the squares vanish, leaving the displayed differential.
Squarefree coordinate monomials and exterior chord monomials give
$2^{n-1}2^\beta=2^m$ basis elements.
Every tree reduction starts from the same completed full edge
complex, giving the asserted zigzag.
\end{proof}

A change of tree need not strictly identify the two finite models.
Symmetries act first on the full edge complex.
A symmetry preserving the chosen tree acts directly on its quotient,
with the induced exterior signs.

\begin{proposition}\label{ei:prop:trace}
Fix orders of the tree coordinates and chord generators.
The top-coefficient functional
\[
 \tau(a)=[u_1\cdots u_{n-1}\epsilon_1\cdots\epsilon_\beta]a
\]
gives a perfect cochain pairing
$C_{G,T}\otimes C_{G,T}\to\mathbb C[\beta]$ by multiplication.
Here $C[r]^i=C^{i+r}$. Consequently
\[
 \dim H^{-j}(C_{G,T})
  =\dim H^{-(\beta-j)}(C_{G,T}).
\]
\end{proposition}

\begin{proof}
Each basis monomial pairs only with its complementary monomial.
The exterior ordering supplies its sign, so the matrix is a signed
permutation matrix. The differential lowers exterior length,
which gives $\tau(da)=0$.
Applying this identity to $d(ab)$ proves compatibility with the
differential. A perfect pairing of finite complexes over a field
induces a perfect pairing on cohomology.
\end{proof}

The normalization of this trace is part of the construction.
A graph symmetry may change it by a character.
We first compute the simplest cycle, then determine this character
in a repeated-color example.

\section{The triangle and its minimal algebra}
\label{ei:gio:sec:operations}

For the triangle, \eqref{ei:eq:triangle-koszul} reduces
along the tree equations $x^2,y^2$ to
\begin{equation}\label{ei:gio:eq:triangle}
 A=\mathbb C[x,y]/(x^2,y^2)\otimes\Lambda(\eta),
 \qquad |x|=|y|=0,\quad |\eta|=-1,\quad d\eta=xy.
\end{equation}
Indeed the remaining equation is $2xy$. Dividing its
odd generator by two gives $\eta$.
The algebra has eight basis elements. Write $H=H^*(A)$, represented by
\begin{equation}\label{ei:gio:eq:triangle-cohomology}
 H^0=\langle1,x,y\rangle,\qquad
 H^{-1}=\langle x\eta,y\eta,xy\eta\rangle.
\end{equation}

The radical $V=\langle x,y\rangle\subset H^0$ has square
zero. Its only nonzero actions on the displayed negative
classes are
\[
 x\cdot y\eta=y\cdot x\eta=xy\eta.
\]
Products of two negative classes vanish.

Let $i:H\to A$ include these representatives, and let
$p:A\to H$ kill $xy$ and $\eta$ and retain the other six
basis elements. Set
\[
 \mathsf h(xy)=\eta,\qquad
 \mathsf h=0\quad\hbox{on the other basis elements}.
\]
Then $pi=1$, $d\mathsf h+\mathsf h d=1-ip$, and
$\mathsf h^2=\mathsf h i=p\mathsf h=0$.
Define a symmetric form $q$ on $V$ by
$q(x,y)=q(y,x)=1$ and $q(x,x)=q(y,y)=0$.

Augment $A$ by $x,y,\eta\mapsto0$, and give $H$ the induced
augmentation. The augmentation ideal is the kernel of this map.

We fix the bar signs before giving the higher product.
Write $sa$ for suspension, with $|sa|=|a|-1$.
On the tensor coalgebra of suspended augmentation ideals
use the degree-one components
\begin{align}
 b_1(sa)&=-s(da),\notag\\
 b_2(sa,sb)&=(-1)^{|a|}s(ab),\notag\\
 b_3(sa,sb,sc)&=(-1)^{2|a|+|b|}s\,m_3(a,b,c).
 \label{ei:gio:eq:bar-signs}
\end{align}
Their extension as coderivations uses the Koszul sign
for passing a degree-one map across the preceding
suspended inputs. The source algebra $A$ has $b_3=0$.
The minimal algebra below has $b_1=0$.

\begin{theorem}\label{ei:gio:thm:triangle-minimal}
The binary product in \eqref{ei:gio:eq:triangle-cohomology}
extends to a strictly unital minimal $A_\infty$ algebra by
\begin{equation}\label{ei:gio:eq:triangle-m3}
 m_3(a,b,c)=\bigl(q(b,c)a-q(a,b)c\bigr)\eta,
 \qquad a,b,c\in V,
\end{equation}
and by setting $m_3=0$ on every other pattern of inputs.
All $m_n$ with $n\ge4$ vanish. A strictly unital
quasi-isomorphism from this algebra to $A$ has suspended
components
\[
 I_1(sa)=s\,i(a),\qquad
 I_2(sa,sb)=s\,\mathsf h(i(a)i(b)),\qquad I_{n\ge3}=0.
\]
The second component is nonzero only for inputs in $V$.
\end{theorem}

\begin{proof}
The binary product is $p(i(a)i(b))$.
The only nonzero homotopy replaces $xy$ by $\eta$.
In arity two the bar comparison equation is exactly
\[
 d\mathsf h(i(a)i(b))=i(a)i(b)-i\,m_2(a,b).
\]

For three inputs in $V$, the two remaining terms are
\[
 p\bigl(i(a)\mathsf h(i(b)i(c))-\mathsf h(i(a)i(b))i(c)\bigr),
\]
which gives \eqref{ei:gio:eq:triangle-m3} with the signs
of \eqref{ei:gio:eq:bar-signs}. An input of degree $-1$ already
contains $\eta$, so both terms vanish.

In every higher transfer tree, a further homotopy above
an $\eta$-marked expression is zero, or two homotopy
outputs are multiplied and give $\eta^2=0$.
Thus all higher products vanish. The same two facts
kill all comparison equations above arity three and
all composites of two ternary operations.

In arity four,
for $a,b,c,d\in V$ the remaining equation is
\[
 m_2(m_3(a,b,c),d)+m_2(a,m_3(b,c,d))=0.
\]
The two coefficients of $xy\eta$ cancel by
\eqref{ei:gio:eq:triangle-m3}. If an input has degree
$-1$, every term vanishes. These formulas give the
transfer along the stated contraction and verify
its bar-square equations. Every higher component
vanishes on a unit, while $I_1(1)=1$.
\end{proof}

The ternary product is detected without choosing a
minimal-model basis. Since $xy=d\eta$ and $y^2=0$,
the triple Massey set is
\begin{equation}\label{ei:gio:eq:triangle-massey}
 \langle[x],[y],[y]\rangle
 =-[y\eta]+\langle[xy\eta]\rangle.
\end{equation}
Its indeterminacy is
$xH^{-1}+H^{-1}y=\langle xy\eta\rangle$.
It does not contain zero, so $A$ is nonformal.
This operation belongs to the simultaneous derived
intersection of the doubled pair divisors. It does not
assert that the original Hall multiplication fails
associativity.

\begin{remark}\label{ei:gio:rem:linear-and-euler}
The multiplicity of the equations matters. With linear
pair equations the corresponding formal complex is
$K(x,y,x+y)$. Subtracting the first two odd generators
from the third makes the latter closed. Resolving $x,y$
then gives $\mathbb C\otimes\Lambda(\epsilon)$ with zero
differential. The closed generator records the linear dependence among the
three equations. The ternary operation in
\eqref{ei:gio:eq:triangle-m3} records the doubled equations.

There is also an exact algebraic comparison with an Euler
extension. Put $X=\mathbb P^1\times\mathbb P^1$ and choose
the two degree-two generators $x,y$ with $\int_Xxy=1$.
Then $H^*(X;\mathbb C)=\mathbb C[x,y]/(x^2,y^2)$ and
$e(TX)=4xy$. In the algebra
$H^*(X;\mathbb C)\otimes\Lambda(e)$ with $|e|=3$ and
$de=4xy$, set $\eta=e/4$.

Its multiplication and
differential are exactly those of \eqref{ei:gio:eq:triangle},
after changing the integer degrees of $x,y,\eta$ to
$0,0,-1$. Parity is unchanged. This comparison of
algebras does not identify their geometric sources.
\end{remark}

\section{Two equal colors}

Take four labeled vertices $0,1,2,3$ with colors
$(1,1,2,3)$, and join precisely the pairs of unequal
colors. The tree with edges $02,12,23$ is preserved
by the transposition $\sigma=(01)$. In logarithmic
coordinates put
\[
 x=z_0-z_2,\qquad y=z_1-z_2,\qquad z=z_3-z_2.
\]
The two chords give $(x-z)^2$ and $(y-z)^2$.
After the tree quotient and rescaling their generators,
the finite model is
\begin{equation}\label{ei:gio:eq:repeated-color}
 C=\mathbb C[x,y,z]/(x^2,y^2,z^2)
           \otimes\Lambda(\eta_1,\eta_2),
 \qquad d\eta_1=xz,\quad d\eta_2=yz.
\end{equation}

All three coordinates have degree zero and both odd
generators degree $-1$. The actual graph transposition
acts by
\[
 \sigma(x)=y,\quad\sigma(y)=x,\quad\sigma(z)=z,
 \qquad\sigma(\eta_1)=\eta_2,
       \sigma(\eta_2)=\eta_1.
\]

In particular $\sigma(\eta_1\eta_2)=-\eta_1\eta_2$.
This action commutes with $d$ and is the action induced
from the full edge complex. Write $C^+$ and $C^-$ for
its invariant and anti-invariant subcomplexes.

\begin{theorem}\label{ei:gio:thm:color-cohomology}
The cohomology dimensions, in degrees $0,-1,-2$, are
\begin{equation}\label{ei:gio:eq:character-dimensions}
 \dim H^*(C)=(5,10,5),\qquad
 \dim H^*(C^+)=(4,5,1),\qquad
 \dim H^*(C^-)=(1,5,4).
\end{equation}
The invariant cohomology algebra has bases
\begin{align*}
 H^0(C^+)&=\langle1,\alpha,z,\chi\rangle,
             &\alpha&=x+y,&\chi&=xy,\\
 H^{-1}(C^+)&=\langle p,q,r,s,t\rangle,\\
 p&=x\eta_1+y\eta_2,&
 q&=z(\eta_1+\eta_2),\\
 r&=xy(\eta_1+\eta_2),&
 s&=xz\eta_1+yz\eta_2,\\
 t&=yz\eta_1+xz\eta_2,\\
 H^{-2}(C^+)&=\langle u\rangle,
             &u&=z(x-y)\eta_1\eta_2.
\end{align*}
Apart from the unit and graded commutativity, its
nonzero products of these basis elements are
\begin{equation}\label{ei:gio:eq:color-products}
 \alpha^2=2\chi,\qquad
 \alpha p=r,\qquad zp=s,\qquad
 \alpha q=s+t,\qquad pq=u.
\end{equation}
\end{theorem}

\begin{proof}
The dimensions of $C$ in degrees $0,-1,-2$ are
$8,16,8$. The image of $d:C^{-1}\to C^0$ is
$\langle xz,yz,xyz\rangle$.
The other differential sends $a\eta_1\eta_2$ to
$a(xz\eta_2-yz\eta_1)$. Its image has basis
\[
 xz\eta_2-yz\eta_1,\quad xyz\eta_1,\quad xyz\eta_2.
\]
Both differentials have rank three, proving the total
dimensions.

The invariant degree-zero space has basis
$1,\alpha,z,\chi,\alpha z,xyz$, and its boundaries
span $\alpha z,xyz$. The invariant cycles of degree
$-1$ have basis $p,q,r,s,t,w$, where
$w=xyz(\eta_1+\eta_2)$. The invariant degree-$-2$
space has basis $(x-y)\eta_1\eta_2$ and $u$.
Their differentials are $-w$ and zero. This proves
the displayed invariant bases.

Since taking invariants
is exact over $\mathbb C$, subtraction gives the
anti-invariant dimensions. Direct multiplication gives
\eqref{ei:gio:eq:color-products}. Every omitted product
contains a squared coordinate or is a multiple of the
boundary $w$.
\end{proof}

\begin{corollary}\label{ei:gio:cor:color-massey}
The invariant algebra $C^+$ is nonformal. More precisely,
\begin{equation}\label{ei:gio:eq:color-massey}
 \langle[\alpha],[z],[z]\rangle
 =-[q]+\langle[r],[s],[t]\rangle
\end{equation}
does not contain zero.
\end{corollary}

\begin{proof}
The invariant primitive $\eta_1+\eta_2$ has differential
$\alpha z$, and $z^2=0$. It gives the representative
$-z(\eta_1+\eta_2)=-q$. By
\eqref{ei:gio:eq:color-products}, the indeterminacy
$\alpha H^{-1}(C^+)+H^{-1}(C^+)z$ is
$\langle r,s,t\rangle$. The cohomology basis shows
that it does not contain $q$.
\end{proof}

\begin{proposition}\label{ei:gio:prop:trace-character}
The perfect trace of this labeled model,
$\tau(a)=[xyz\eta_1\eta_2]a$, satisfies
$\tau(\sigma a)=-\tau(a)$. Its pairing restricts to
perfect pairings between opposite characters, and
vanishes between equal characters. The invariant
cohomology has no perfect scalar Frobenius pairing
supported in total degree $-2$.
\end{proposition}

\begin{proof}
The coordinate product $xyz$ is fixed and the exterior
product changes sign. Thus
$\tau(\sigma a\,\sigma b)=-\tau(ab)$.
If $a,b$ have characters $\varepsilon,\delta$, a
nonzero pairing requires $\varepsilon\delta=-1$.
The full pairing is perfect by
Proposition~\ref{ei:prop:trace}. Decomposition into
the two characters proves perfectness of the opposite
blocks. Such a scalar Frobenius trace on the
invariant cohomology would pair $H^0$ perfectly with
$H^{-2}$, but these have dimensions four and one.
\end{proof}

The character is part of this explicit descent problem.
It does not rule out a comparison retaining a nontrivial
orientation line or both character sectors. It rules
out discarding them while keeping the same perfect
scalar trace.

\section{Every graph with one cycle}

The triangle dimensions extend to all connected graphs
with one cycle. Suppose that the cycle has length $c$
and that the graph has $n$ vertices. Removing a cycle
edge gives a spanning tree. Orient its $c-1$ edges
along the remaining cycle path and put
\[
 A_s=\mathbb C[u_1,\ldots,u_s]/(u_1^2,\ldots,u_s^2),
 \qquad s=c-1,\qquad L=u_1+\cdots+u_s.
\]
The one chord has equation $L^2$. The $n-c$ tree
coordinates outside the cycle occur in no chord equation.

For polynomial degree, give every $u_a$ degree one.
This grading differs from cochain degree.

\begin{theorem}\label{ei:gio:thm:unicyclic}
For this graph, cohomology occurs only in degrees
$0,-1$, and
\begin{equation}\label{ei:gio:eq:cycle-dimensions}
 \dim H^0(C_{G,T})=\dim H^{-1}(C_{G,T})
 =2^{n-c}\binom{c}{\lfloor c/2\rfloor}.
\end{equation}
For the cycle alone, the polynomial-degree-$k$ part
of $H^0$ has dimension
\[
 \max\left\{\binom{s}{k}-\binom{s}{k-2},0\right\},
\]
where a binomial coefficient outside $0,\ldots,s$ is zero.
\end{theorem}

\begin{proof}
On the squarefree monomials $u_I$ define
$F(u_I)=\sum_{i\in I}u_{I\setminus\{i\}}$.
Multiplication by $L$ and this operator satisfy
\[
 [F,L](u_I)=(s-2|I|)u_I.
\]
With $H(u_I)=(s-2|I|)u_I$, one also has
$[H,L]=-2L$ and $[H,F]=2F$.
These identities can be checked one variable at a
time. The algebra $A_s$ is the tensor product of the resulting
two-dimensional representations.

For completeness, give its monomial basis the
orthonormal positive Hermitian form. Then $F=L^*$.
A primitive vector $v$ of degree $j$, meaning $Fv=0$,
generates a string satisfying
\[
 F L^r v=r(s-2j-r+1)L^{r-1}v.
\]
The formula follows inductively from the commutators.

Its norm identity shows that the nonzero string
runs from degree $j$ to degree $s-j$. Its span is
stable under $F,L,H$, so its orthogonal complement
is stable too. Repetition decomposes $A_s$ into such
symmetric strings.

Consequently $L^2:(A_s)^{k-2}\to(A_s)^k$ has maximal rank.
On a string it is nonzero exactly when both positions occur.
The sets of strings containing these two positions are nested. This proves the cokernel formula.

Summing its positive differences telescopes to
\[
 \binom{s}{\lfloor s/2\rfloor}
 +\binom{s}{\lfloor(s-1)/2\rfloor}
 =\binom{s+1}{\lfloor(s+1)/2\rfloor}.
\]
The kernel and cokernel of the endomorphism $L^2$
have the same dimension. The remaining $n-c$
squarefree variables contribute the factor $2^{n-c}$.
\end{proof}

Thus a tree on four vertices has only $H^0$ of
dimension eight, a triangle has dimensions $(3,3)$,
and a four-cycle has $(6,6)$. The repeated-color
graph above has $(5,10,5)$.

For comparison take the complete graph on four
distinct colors. Its star-tree model is
\[
 A_3\otimes\Lambda(\alpha,\beta,\gamma),\qquad
 d\alpha=xy,\quad d\beta=xz,\quad d\gamma=yz,
\]
where $A_3=\mathbb C[x,y,z]/(x^2,y^2,z^2)$.
The three differential ranks, from degrees $-1,-2,-3$,
are $4,6,4$. The first image is
$\langle xy,xz,yz,xyz\rangle$.

For the middle image use
\[
 xy\beta-xz\alpha,\quad xy\gamma-yz\alpha,
 \quad xz\gamma-yz\beta,
 \quad xyz\alpha,\quad xyz\beta,\quad xyz\gamma.
\]
These six independent vectors span all products of
the three basic images. 
The last kernel is the
common annihilator of $xy,xz,yz$, namely
$\langle xy,xz,yz,xyz\rangle$.
Since the term dimensions are $8,24,24,8$, the
cohomology dimensions are $(4,14,14,4)$.
This is a four-color example. All these numbers
measure the specified incidence coefficients,
rather than root multiplicities of an Igusa algebra.

\section{Independent quadratic parameters}
\label{ei:sec:independent}

The finite graph complex can occur as a derived fiber of a family
whose total cohomology is concentrated in degree zero. We first give
each edge its own quadratic parameter. The actual elliptic
displacement will then provide a different comparison.


Retain the connected graph $G$, spanning tree $T$,
tree coordinates $u_a$ and chord forms $\ell_c$ of
Theorem~\ref{ei:thm:finite-model}.
For a tree edge $a$ also write $\ell_a=u_a$.
Give every edge an independent parameter $t_e$, and put
\[
 R_G=\mathbb C[t_e:e\in E(G)],\qquad
 B_T=R_G[u_1,\ldots,u_{n-1}]/(u_a^2-t_a:a\in T).
\]
The squarefree monomials $u_I$ form an $R_G$-basis of
$B_T$. Define
\begin{equation}\label{ei:gif:eq:graph-family}
 \mathcal C_{G,T}=B_T\otimes\Lambda(\epsilon_c:c\notin T),
 \qquad d\epsilon_c=\ell_c^2-t_c,
 \qquad |\epsilon_c|=-1.
\end{equation}

This is a polynomial deformation of the finite formal
model. The parameters are independently assigned to
the edges. No identification with a universal compact
Hall deformation is part of its definition.

\begin{theorem}\label{ei:gif:thm:graph-family}
The complex \eqref{ei:gif:eq:graph-family} is finite free
over $R_G$, with term ranks
$2^{n-1}\binom{\beta}{k}$ in degree $-k$.
Its only cohomology is
\begin{equation}\label{ei:gif:eq:family-cohomology}
 H^0(\mathcal C_{G,T})\cong
 \mathbb C[u_1,\ldots,u_{n-1}],\qquad
 t_e\longmapsto\ell_e^2.
\end{equation}
Its derived fiber at $t_e=0$ is the graph algebra
$C_{G,T}$. The coefficient functional
\[
 a\longmapsto
 [u_1\cdots u_{n-1}\epsilon_1\cdots\epsilon_\beta]a
\]
gives a perfect pairing
$\mathcal C_{G,T}\otimes_{R_G}\mathcal C_{G,T}
\to R_G[\beta]$ of complexes.
\end{theorem}

\begin{proof}
Eliminate the tree parameters to identify
$B_T$ with the polynomial ring
$\mathbb C[u_1,\ldots,u_{n-1},t_c:c\notin T]$.
Each remaining equation $\ell_c^2-t_c$ eliminates
its own parameter. In any order they form a regular
sequence, with quotient the ring in
\eqref{ei:gif:eq:family-cohomology}. Its Koszul resolution
is precisely \eqref{ei:gif:eq:graph-family}.
The displayed free basis gives the term ranks and
computes derived base change by ordinary tensoring.
At the origin its differential reduces to
\eqref{ei:eq:finite-model}.

For the pairing, a product $u_Iu_J$ contributes to
the top squarefree coefficient exactly when
$J=I^c$. An overlap replaces a square by a parameter
and loses that coordinate from the squarefree part.
The exterior factors likewise pair complementary
subsets. The pairing matrix is therefore a signed
permutation matrix over $R_G$. Its trace vanishes on
the differential because exterior length decreases.
The derivation rule proves compatibility with $d$.
\end{proof}

The theorem also holds after completion at the parameter
origin. The image ideal contains every $u_a^2$, so
the completed degree-zero cohomology is
$\mathbb C[[u_1,\ldots,u_{n-1}]]$.
The family need not be flat in cohomology: its central
derived fiber retains the negative classes computed
in Section~\ref{ei:gio:sec:operations}.
On a cover adjoining square roots of the parameters,
the equations factor as
$(\ell_e-\sqrt{t_e})(\ell_e+\sqrt{t_e})$.
No such cover is needed to define the family itself.

\subsection{The triangle discriminant}

For three edges set
$t_{12}=a$, $t_{23}=b$, and $t_{13}=a+b+2c$.
After rescaling the chord generator, the family is
\begin{equation}\label{ei:gif:eq:triangle-family}
 R=\mathbb C[a,b,c],\qquad
 B=R[x,y]/(x^2-a,y^2-b),\qquad
 \mathcal A_R=B\otimes\Lambda(\eta),\quad d\eta=xy-c.
\end{equation}
Thus $\mathcal A_R$ is free of rank four in each of
degrees $-1,0$. Its total cohomology is
$\mathbb C[x,y]$, with $a=x^2$, $b=y^2$, $c=xy$.
In the basis $1,x,y,xy$, multiplication by $xy-c$ is
\begin{equation}\label{ei:gif:eq:triangle-matrix}
 F=\begin{pmatrix}
 -c&0&0&ab\\
 0&-c&b&0\\
 0&a&-c&0\\
 1&0&0&-c
 \end{pmatrix},\qquad
 \det F=(c^2-ab)^2.
\end{equation}
Put $\delta=c^2-ab$ and $S=R/(\delta)$.

\begin{theorem}\label{ei:gif:thm:triangle-module}
As an $S$-module,
\begin{equation}\label{ei:gif:eq:nonfree-module}
 H^0(\mathcal A_R)\cong S\oplus I,
 \qquad I=(a,c)\subset S.
\end{equation}
The second summand has rank one and is nonfree at the
cone vertex. It has the two-periodic free resolution
over $S$ determined by
\begin{equation}\label{ei:gif:eq:matrix-factorization}
 M=\begin{pmatrix}-c&b\\a&-c\end{pmatrix},\qquad
 N=\begin{pmatrix}-c&-b\\-a&-c\end{pmatrix},
 \qquad MN=NM=\delta I_2.
\end{equation}
The map $\operatorname{Spec}\mathbb C[x,y]
\to\operatorname{Spec}S$ is the quotient double cover
by $(x,y)\mapsto(-x,-y)$. The ring $S$ is already normal.
\end{theorem}

\begin{proof}
Split the four basis elements by this involution.
On $\langle1,xy\rangle$, the matrix is
$\left(\begin{smallmatrix}-c&ab\\1&-c\end{smallmatrix}\right)$.
Its unit entry permits invertible row and column
operations reducing it to $\operatorname{diag}(1,-\delta)$.
On $\langle x,y\rangle$, the matrix is $M$.
The invariant subring of $\mathbb C[x,y]$ is generated
by $x^2,y^2,xy$, with relation $c^2=ab$. It is $S$.

Multiplication by $x$ identifies the anti-invariant
summand with $(a,c)$, sending $x$ to $a$ and $y$ to $c$.
This multiplication is injective in $\mathbb C[x,y]$.
The odd block of \eqref{ei:gif:eq:triangle-matrix} supplies
its presentation by $M$.

The ideal $I$ is nonzero in the domain $S$, so it has
rank one. At the maximal ideal $(a,b,c)$ its classes
$a,c$ are independent modulo $(a,b,c)I$. It therefore
requires two generators there and cannot be a free
rank-one module.

The products in \eqref{ei:gif:eq:matrix-factorization}
follow by multiplication. To prove exactness, suppose
$M\bar v=0$ over $S$. Lift to $R$, so
$Mv=\delta w=MNw$. The determinant of $M$ is the
nonzero polynomial $\delta$, hence $M$ is injective
over $R$ and $v=Nw$. Thus $\ker\bar M=\operatorname{im}\bar N$.
Interchanging $M,N$ proves the other equality.

Finally, let an element of the invariant fraction field be
integral over the invariant ring. It is also integral over the
normal ring $\mathbb C[x,y]$, so it belongs to that ring.
It is still invariant, proving that the invariant ring is normal. The invariant
field has index two in $\mathbb C(x,y)$.
The cover is consequently not a normalization map.
\end{proof}

In this decomposition, invariant and anti-invariant
refer to the simultaneous sign action, not to cochain
degree or super parity. A generic rank-two module
has here a nonfree rank-one summand. Replacing it by
two free copies would lose the module at the vertex.

\begin{proposition}\label{ei:gif:prop:triangle-fibers}
The derived fibers of \eqref{ei:gif:eq:triangle-family}
have the following dimensions:
\[
\begin{array}{c|ccc}
 \text{parameter}&\operatorname{rank}F&\dim H^0&\dim H^{-1}\\\hline
 \delta\ne0&4&0&0\\
 \delta=0,\ (a,b,c)\ne0&2&2&2\\
 (a,b,c)=0&1&3&3.
\end{array}
\]
Every smooth-discriminant fiber is formal, with
cohomology algebra $\mathbb C^2\otimes\Lambda(\xi)$,
$|\xi|=-1$. The vertex fiber is the nonformal
triangle algebra \eqref{ei:gio:eq:triangle}.
\end{proposition}

\begin{proof}
Off the discriminant, $F$ is invertible.
At a nonzero point of $\delta=0$, at least one of
$a,b$ is nonzero. If $a\ne0$, split $x^2=a$ into
the two roots $r,-r$. Each component is
\[
 \mathbb C[y]/(y^2-b)\otimes\Lambda(\eta),
 \qquad d\eta=r_jy-c,\qquad r_j\in\{r,-r\}.
\]
Multiplication by $r_jy-c$ has rank one.

Its kernel is the line generated by $r_jy+c$,
since $r_j^2b=c^2$ and the coefficient $r_j$ is
nonzero. Sending $\xi$ to $(r_jy+c)\eta$ and
$1$ to $1$ gives a differential graded algebra
quasi-isomorphism from $\Lambda(\xi)$ to this
component. The case $b\ne0$ is symmetric.
At the vertex, $F$ has only its bottom-left entry
nonzero. Its fiber is exactly
\eqref{ei:gio:eq:triangle}, whose Massey product was
computed in \eqref{ei:gio:eq:triangle-massey}.
\end{proof}

Ordinary specialization of $H^0(\mathcal A_R)$ at
the vertex gives only
$\mathbb C[x,y]/(x^2,y^2,xy)$. It loses both the
negative classes and their ternary operation.
The relative pairing can also be seen directly.
Extraction of the coefficient of $xy$ in $B$ has
Gram matrix
\[
 P=\begin{pmatrix}
 0&0&0&1\\0&0&1&0\\0&1&0&0\\1&0&0&0
 \end{pmatrix},\qquad F^tP=PF.
\]
Extracting additionally $\eta$ gives the perfect
pairing $\mathcal A_R\otimes_R\mathcal A_R\to R[1]$.
The jump in fiber cohomology therefore occurs inside
a fixed free complex with a perfect chain pairing.

\subsection{A nonnormal four-cycle image}

For a four-cycle choose consecutive tree coordinates
$u,v,w$. Its final edge has difference $u+v+w$.
Write its parameter as $a+b+c+2z$, where the three
tree parameters are $a,b,c$. Its finite free family is
\begin{equation}\label{ei:gif:eq:four-cycle-family}
 \begin{gathered}
 R_4=\mathbb C[a,b,c,z],\qquad
 B_4=R_4[u,v,w]/(u^2-a,v^2-b,w^2-c),\\
 \mathcal A_4=B_4\otimes\Lambda(\eta),\qquad
 |\eta|=-1,\qquad d\eta=uv+uw+vw-z.
 \end{gathered}
\end{equation}
It resolves $\mathbb C[u,v,w]$ with parameter map
\[
 (a,b,c,z)=(u^2,v^2,w^2,uv+uw+vw).
\]
Define the quartic polynomial
\begin{equation}\label{ei:gif:eq:quartic}
 P_4=(z^2+ab-c(a+b))^2-4ab(z+c)^2.
\end{equation}

\begin{theorem}\label{ei:gif:thm:four-cycle-normalization}
The parameter image in \eqref{ei:gif:eq:four-cycle-family}
has coordinate ring
\[
 A_4=\mathbb C[a,b,c,z]/(P_4),
\]
and $P_4$ is irreducible. If $F_4$ is multiplication
by $uv+uw+vw-z$ on the rank-eight $R_4$-module $B_4$,
then $\det F_4=P_4^2$.
The ring $A_4$ is not normal. Its normalization is
\begin{equation}\label{ei:gif:eq:normalization}
 D_4=\mathbb C[u,v,w]^{\{\pm1\}},
\end{equation}
where the sign acts simultaneously on all three
variables. The map from $\mathbb A^3$ to
$\operatorname{Spec}D_4$ still has generic degree two.
\end{theorem}

\begin{proof}
Put $z=uv+uw+vw$ and observe that
$(z-uv)^2-c(u+v)^2=0$.
The two polynomials $(z-uv)^2-c(u+v)^2$ and
$(z+uv)^2-c(u-v)^2$ are conjugate under the sign
change of $uv$. Their product, after using
$u^2=a,v^2=b,w^2=c$, is
\[
 \bigl(z^2+ab-c(a+b)-2uv(z+c)\bigr)
 \bigl(z^2+ab-c(a+b)+2uv(z+c)\bigr)=P_4.
\]
Thus the four possible values of $uv+uw+vw$ under
independent signs, modulo simultaneous sign, are
the roots of $P_4$.

Over $\mathbb C(a,b,c)$ the extension obtained by
adjoining $u,v,w$ is multiquadratic of degree eight.
Its sign group fixes $uv+uw+vw$ exactly when all
three signs are equal: the three monomials $uv,uw,vw$
are independent over $\mathbb C(a,b,c)$.
Its orbit therefore has four elements, proving that
$P_4$ is the minimal polynomial. Since $P_4$ is
monic in $z$, polynomial division shows that the
kernel of the parameter map is exactly $(P_4)$.

The determinant of multiplication by
$uv+uw+vw-z$ is its norm. Each of the four orbit
values occurs twice among the eight sign choices,
so this norm is $P_4^2$. The identity over the
fraction field is an identity over $R_4$.

The invariant ring $D_4$ is generated by
$u^2,v^2,w^2,uv,uw,vw$. Each cross product is
integral over $A_4$, for example $(uv)^2=ab$.
They also lie in its fraction field, since
\begin{equation}\label{ei:gif:eq:cross-product}
 uv=\frac{z^2+ab-c(a+b)}{2(z+c)},
\end{equation}
and cyclically for $uw,vw$. The denominators are
nonzero elements of that field. Hence $D_4$ is
finite and birational over $A_4$.
It is normal by the same finite-invariant argument
used for the triangle, and is therefore the
integral closure of $A_4$.

Finally, grade $u,v,w$ by degree one. The degree-two
part of the image ring $A_4$ is spanned by
\[
 u^2,\quad v^2,\quad w^2,\quad uv+uw+vw.
\]
It does not contain $uv$. This integral element is
missing from $A_4$, proving nonnormality. Simultaneous
sign has order two on the fraction field of
$\mathbb C[u,v,w]$, so the last assertion follows.
\end{proof}

For the triangle the squares and the single cross
product already generate the sign-invariant ring.
For the four-cycle the parameter records only the
sum of three cross products. Formula
\eqref{ei:gif:eq:cross-product} recovers them generically,
while normalization restores them across the divisors
where those fractions are not regular expressions.
The normalization and the double cover thus record
different geometric operations.


\section{The actual elliptic displacement}
\label{ei:sec:displacement}

Take three ordered distinct colors with arrows pointing
from lower to higher color. Remove the common spectral
translation and use differences $x=z_2-z_1$, $y=z_3-z_2$.
In a formal logarithm the third difference is $x+y$.
Write $h$ also for the formal displacement parameter and set
$R=\mathbb C[[h]]$. Consider the Koszul dg algebra over $R[x,y]$ of
\begin{equation}\label{ei:esf:eq:triangle-equations}
 x(x+h),\qquad y(y+h),\qquad(x+y)(x+y+h).
\end{equation}

Its extension to $\mathbb C[[h,x,y]]$ is the formal elliptic
incidence complex after removing invertible theta factors.
Each Koszul generator has cohomological degree $-1$ and
differential equal to its equation. The differential is a
graded derivation.

\begin{theorem}\label{ei:esf:thm:flat-triangle}
The Koszul algebra of \eqref{ei:esf:eq:triangle-equations} is
multiplicatively quasi-isomorphic to
\begin{equation}\label{ei:esf:eq:triangle}
 B_h=R[x,y]/(x^2+hx,y^2+hy),\qquad
 A_h=B_h\otimes\Lambda(\eta),\qquad
 |\eta|=-1,\quad d\eta=xy.
\end{equation}
It is free of rank four in degrees $0$ and $-1$.
Its cohomology $H=H^*(A_h)$ is free over $R$, with bases
\begin{equation}\label{ei:esf:eq:cohomology}
 H^0=R\{1,x,y\},\qquad
 H^{-1}=R\{k_1,k_2,k_3\},
 \quad
 \begin{aligned}
 k_1&=[(x+h)\eta],\\
 k_2&=[(y+h)\eta],\\
 k_3&=[(xy-h^2)\eta].
 \end{aligned}
\end{equation}
The degree-zero relations are $x^2=-hx$, $y^2=-hy$, $xy=0$.
The action on negative cohomology is
\begin{equation}\label{ei:esf:eq:action}
\begin{array}{c|ccc}
 &k_1&k_2&k_3\\\hline
 x&0&k_3+hk_1&-hk_3-h^2k_1\\
 y&k_3+hk_2&0&-hk_3-h^2k_2.
\end{array}
\end{equation}
Every product of two negative classes vanishes.
\end{theorem}

\begin{proof}
The first two equations form a regular sequence: each is
monic in its own variable, also after quotienting by the
other. Their Koszul algebra resolves $B_h$. Tensoring with
the remaining free Koszul factor preserves this
quasi-isomorphism. In $B_h$ the third equation is $2xy$.
Dividing its odd generator by two gives $\eta$.
The same reduction applies to the completed coordinate
ring, since the resulting finite $R$-module is complete.

In the ordered basis $1,x,y,xy$ of $B_h$, multiplication by
$xy$ has only one nonzero row, namely
$(1,-h,-h,h^2)$, in the $xy$ row. Its cokernel is free on
$1,x,y$. The elements $x+h$, $y+h$, $xy-h^2$ are killed by
$xy$. Subtracting their multiples from any kernel vector
leaves a scalar killed by the unit entry of the displayed
row. Thus they form a free basis of the kernel, proving
\eqref{ei:esf:eq:cohomology}.

Multiplying these representatives in $B_h\otimes\Lambda(\eta)$
gives \eqref{ei:esf:eq:action}. The relation $\eta^2=0$ gives
the last claim.
\end{proof}

The free cohomology has a contraction defined over $R$.
Let $i:H\to A_h$ include the displayed representatives.
On the basis $1,x,y,xy,\eta,x\eta,y\eta,xy\eta$, let
$p:A_h\to H$ kill $xy,\eta$, send the last three elements
to $k_1,k_2,k_3$, and fix $1,x,y$. Put
\begin{equation}\label{ei:esf:eq:contraction}
 \mathsf H(xy)=\eta,\qquad
 \mathsf H=0\text{ on the other basis elements}.
\end{equation}
Then $pi=1$, $d\mathsf H+\mathsf Hd=1-ip$, and
$\mathsf H^2=\mathsf Hi=p\mathsf H=0$.
For example,
$\mathsf Hd(x\eta)=-h\eta=x\eta-i(k_1)$,
so the parameter term is part of the contraction.

\begin{proposition}\label{ei:esf:prop:trace}
The functional $\tau:A_h\to R[1]$ extracting the coefficient
of $xy\eta$ defines a perfect cyclic pairing of complexes
by $(a,b)\mapsto\tau(ab)$. The induced pairing between
$H^0$ and $H^{-1}$ is perfect over $R$.
\end{proposition}

\begin{proof}
The multiplication-trace matrix on $B_h$, with trace equal
to the coefficient of $xy$, is
\[
 \begin{pmatrix}
 0&0&0&1\\0&0&1&-h\\0&1&0&-h\\1&-h&-h&h^2
 \end{pmatrix}.
\]
Its determinant is $1$. It pairs the degree-zero and
degree-minus-one copies of $B_h$ in $A_h$.
The trace kills $dA_h$, and the derivation rule gives
$\tau((da)b)+(-1)^{|a|}\tau(a\,db)=0$.
Graded commutativity gives cyclicity. On the cohomology
bases of \eqref{ei:esf:eq:cohomology}, the pairing matrix is
\[
 \begin{pmatrix}0&0&1\\0&1&-h\\1&0&-h\end{pmatrix},
\]
of determinant $-1$.
\end{proof}

\subsection{The ternary class at coalescence}

The contraction determines the multiplication beyond
cohomology. Put $V=Rx\oplus Ry\subset H^0$, and define
$q:V\otimes_R V\to R$ by
$q(x,y)=q(y,x)=1$, $q(x,x)=q(y,y)=0$.
Let $\psi:V\to H^{-1}$ send $x$ to $k_1$ and $y$ to $k_2$.

Augment $A_h$ by $x,y,\eta\mapsto0$ and $H$ by
$x,y,k_1,k_2,k_3\mapsto0$.
For $A=A_h$ or $H$, put $\overline A=\ker(A\to R)$.
On $T^c_R(s\overline A)=\bigoplus_{n\ge1}
(s\overline A)^{\otimes_Rn}$ use the suspension
$|sa|=|a|-1$ and bar components
\[
 b_1(sa)=-s(da),\quad
 b_2(sa,sb)=(-1)^{|a|}s(ab),\quad
 b_3(sa,sb,sc)=(-1)^{2|a|+|b|}s\,m_3(a,b,c).
\]
Their coderivation extensions use the Koszul sign for
passing a degree-one map across the preceding suspended
inputs. For $A_h$ set $b_3=0$, and for $H$ set $b_1=0$.

\begin{theorem}\label{ei:esf:thm:minimal}
The multiplication \eqref{ei:esf:eq:action} extends to a
strictly unital minimal $A_\infty$ algebra on $H$ by
\begin{equation}\label{ei:esf:eq:ternary}
 m_3(a,b,c)=q(b,c)\psi(a)-q(a,b)\psi(c),\qquad a,b,c\in V.
\end{equation}
The ternary product is zero on all other basis input
patterns, and $m_n=0$ for $n\ge4$.
A strictly unital quasi-isomorphism to $A_h$ has suspended
components
\[
 I_1(sa)=s\,i(a),\qquad
 I_2(sa,sb)=s\,\mathsf H(i(a)i(b)),\qquad I_{n\ge3}=0.
\]
\end{theorem}

\begin{proof}
The binary transfer is $p(i(a)i(b))$.
In arity two the comparison equation is
$d\mathsf H(i(a)i(b))=i(a)i(b)-i\,m_2(a,b)$.
For three degree-zero inputs, the two transfer trees give
\[
 p\bigl(i(a)\mathsf H(i(b)i(c))
       -\mathsf H(i(a)i(b))i(c)\bigr).
\]
Since $\mathsf H$ extracts the $xy$ coefficient and
$p(x\eta)=k_1$, $p(y\eta)=k_2$, this is
\eqref{ei:esf:eq:ternary}. A negative input makes both terms
zero. Unit inputs vanish because $p\mathsf H=\mathsf Hi=0$.

Two internal homotopies in a transfer tree either lie on
the same branch or on different branches. On the same
branch the second acts on an expression containing $\eta$
and vanishes. On different branches their outputs
multiply to zero because $\eta^2=0$.
All higher transfer products therefore vanish.

For the comparison, the final homotopy
also kills each expression containing $\eta$, so
$I_{n\ge3}=0$. The bar comparison in arity three is the
displayed two-tree formula together with
$d\mathsf H+\mathsf Hd=1-ip$. In every higher arity its
terms contain either two $I_2$ outputs or a homotopy
applied to a negative input and vanish. Thus the bar
comparison holds in every arity.

Its leading tensor-length
term is $i^{\otimes n}$, so the induced coalgebra map is
injective. Compatibility with the square-zero
bar differential on $A_h$ implies $b_H^2=0$.
Finally $I_1$ induces the identity on cohomology, and
every higher component vanishes on a unit.
\end{proof}

The ternary operation has a more precise relative invariant
than its nonvanishing formula. Use the normalized Hochschild
complex of the graded $R$-algebra $(H,m_2)$. An arity-two
cochain of internal degree $-1$ can be nonzero only on
$V\otimes V$, with values in $H^{-1}$. Define such a cochain by
\begin{equation}\label{ei:esf:eq:primitive}
 G_2(x,x)=0,\quad G_2(x,y)=G_2(y,x)=k_1,\quad G_2(y,y)=k_2.
\end{equation}
For degree-zero inputs its coboundary is
\[
 (\delta_HG_2)(a,b,c)=aG_2(b,c)-G_2(ab,c)
                       +G_2(a,bc)-G_2(a,b)c.
\]

\begin{proposition}\label{ei:esf:prop:torsion-class}
The arity-three, internal-degree-minus-one Hochschild
class of $m_3$ is nonzero and has annihilator exactly $(h)$.
More precisely,
\begin{equation}\label{ei:esf:eq:torsion-class}
 \delta_HG_2=h\,m_3.
\end{equation}
\end{proposition}

\begin{proof}
The arity-four $A_\infty$ identity says that $m_3$ is a
Hochschild cocycle. On the ordered triples
$xxx,xxy,xyx,xyy,yxx,yxy,yyx,yyy$, the coboundary of
\eqref{ei:esf:eq:primitive} is
\[
 h(0,k_1,0,-k_2,-k_1,0,k_2,0),
\]
by \eqref{ei:esf:eq:action}. These are exactly the values
of $hm_3$. Other normalized input patterns vanish by
degree, proving \eqref{ei:esf:eq:torsion-class}.

Reduce modulo $h$, and let $G$ be a normalized two-cochain
of internal degree $-1$. Since $xy=y^2=0$,
\[
 (\delta_HG)(x,y,y)=xG(y,y)-G(x,y)y\in\mathbb Ck_3:
\]
the action of $x,y$ on negative classes has image in that line.
But $m_3(x,y,y)=-k_2$. Hence $m_3$ is not a boundary
after reduction, and therefore is not a boundary over $R$.
Every scalar outside $(h)$ is a unit in $R$.
Equation~\eqref{ei:esf:eq:torsion-class} therefore gives
the exact annihilator.
\end{proof}

Equivalently, at $h=0$ the triple Massey set is
\[
 \langle[x],[y],[y]\rangle=-k_2+\mathbb Ck_3.
\]
Indeed $xy=d\eta$, $y^2=0$, and the indeterminacy is
$xH^{-1}+H^{-1}y=\mathbb Ck_3$. It excludes zero.
Thus the central algebra is nonformal even though the
relative cohomology is free and commutes with base change.

\begin{remark}\label{ei:rem:unary-class}
The torsion assertion concerns this particular ternary class.
It does not extend to all positive Hochschild arities.
Indeed let $D$ act by zero on $H^0$ and by the identity on $H^{-1}$.
It is an $R$-linear derivation of internal degree zero: multiplication
by $H^0$ preserves $H^{-1}$, and two negative classes multiply to zero.

In the Hochschild complex of $(H,m_2)$, $D$ is therefore a unary
cocycle. Unary boundaries of internal degree zero are inner
commutators, which vanish because $H$ is graded commutative.
Since $H^{-1}$ is free and nonzero over $R$, no nonzero scalar
annihilates this class. It remains nonzero after inverting $h$.
\end{remark}

\subsection{Localization and the independent quadratic family}

Let $K=\mathbb C((h))$. The tree algebra $B_h\otimes_R K$
has four orthogonal idempotents
\begin{equation}\label{ei:esf:eq:idempotents}
 \begin{aligned}
 e_{00}&=(x+h)(y+h)/h^2,&e_{10}&=-x(y+h)/h^2,\\
 e_{01}&=-(x+h)y/h^2,&e_{11}&=xy/h^2.
 \end{aligned}
\end{equation}
They sum to $1$ and evaluate at the four points with
$x,y\in\{0,-h\}$. On the first three, $xy=0$, and on
the last, $xy=h^2$.

\begin{proposition}\label{ei:esf:prop:generic}
Projection to the first three idempotent factors gives
a multiplicative quasi-isomorphism
\[
 A_h\otimes_R K\longrightarrow
 \prod_{\substack{(i,j)\in\{0,1\}^2\\(i,j)\ne(1,1)}}
 K\otimes\Lambda(\eta_{ij}),
 \qquad |\eta_{ij}|=-1,\quad d=0.
\]
The kernel contracts by multiplication by
$h^{-2}e_{11}\eta$.
\end{proposition}

\begin{proof}
On the $e_{11}$ factor,
$d(h^{-2}e_{11}\eta)=e_{11}$, so the derivation rule
gives a contracting homotopy. The other factors have
zero differential. The product decomposition is by
central idempotents, hence its projection is multiplicative.
\end{proof}

The poles in \eqref{ei:esf:eq:idempotents} prevent this
decomposition from extending across $h=0$. There is also
an explicit logarithmic connection on the punctured model:
\begin{equation}\label{ei:esf:eq:connection}
 \nabla_{\partial_h}
 =\partial_h+h^{-1}
       (x\partial_x+y\partial_y+2\eta\partial_\eta).
\end{equation}
It preserves $x^2+hx,y^2+hy$ and commutes with $d\eta=xy$.
The elements $x/h,y/h,\eta/h^2$ are horizontal.
The residue on the cohomology basis has weights
$0,1,1,3,3,4$. In the analytic realization on a punctured
disk its monodromy is trivial.

This specifies an Euler
connection. An identification with a Gauss--Manin
connection would require an additional comparison.

The equations $X^2=a$, $Y^2=b$ with chord differential $XY-c$,
for independent complex parameters $a,b,c$, give a different family.
The discrepancy already appears under recentering:
with $X=x+h/2$, $Y=y+h/2$ one has
\[
 x^2+hx=X^2-h^2/4,\quad y^2+hy=Y^2-h^2/4,\quad
 xy=XY-\frac h2(X+Y)+\frac{h^2}{4}.
\]
The linear term in the last equation is indispensable.
There is in fact an obstruction to every base map, not
only to this attempted change of coordinates.

\begin{proposition}\label{ei:esf:prop:no-quadratic-pullback}
No map $\mathbb C[a,b,c]\to R$ makes the derived pullback
of
\[
 \mathbb C[a,b,c][X,Y]/(X^2-a,Y^2-b)
       \otimes\Lambda(\zeta),\qquad |\zeta|=-1,\quad d\zeta=XY-c,
\]
quasi-isomorphic to $A_h$ as a dg $R$-algebra.
\end{proposition}

\begin{proof}
Over $K$, the differential of the pulled-back complex,
in the basis $1,X,Y,XY$, is
\[
 F=\begin{pmatrix}
 -c&0&0&ab\\0&-c&b&0\\0&a&-c&0\\1&0&0&-c
 \end{pmatrix}.
\]
Both cohomology dimensions can be $3$ only if $F$ has
rank $1$. Its two-by-two minors include $a,b,c$ up to
sign, so this forces $a=b=c=0$ in $K$ and therefore in $R$.
The pullback would then be the constant central algebra.
Its Massey set $-k_2+Kk_3$ still excludes zero, so it
remains nonformal over $K$. Proposition~\ref{ei:esf:prop:generic}
makes the generic fiber of $A_h$ formal, a contradiction.
\end{proof}


\section{Displacement along signed graph paths}
\label{ei:sec:shifted-graphs}

The triangle reduction extends to every connected finite
simple graph $G$ on ordered vertices $0,\ldots,n-1$,
with $R=\mathbb C[[h]]$ as in the formal triangle.
Orient each edge from its lower to its higher vertex,
and assign it the pair equation $\ell_e(\ell_e+h)$,
where $\ell_{ij}=z_j-z_i$. Set $z_0=0$ to remove common
spectral translation, choose an oriented spanning tree
$T$, and use its $k=n-1$ edge differences $u_1,\ldots,u_k$
as coordinates. For a chord $c$ write
\[
 \ell_c=\sum_{a=1}^k\sigma_{ca}u_a,\qquad
 \sigma_{ca}\in\{-1,0,1\},\qquad
 \beta=|E(G)|-n+1.
\]
The signs record the traversal of its tree path.

\begin{theorem}\label{ei:esf:thm:graphs}
The Koszul algebra of the shifted edge equations has a
multiplicative quasi-isomorphism to
\begin{equation}\label{ei:esf:eq:graph-model}
 C_{G,h}=
 R[u_1,\ldots,u_k]/(u_a(u_a+h):1\le a\le k)
       \otimes\Lambda(\epsilon_c:c\notin T),
 \qquad |\epsilon_c|=-1,
\end{equation}
with differential
\begin{equation}\label{ei:esf:eq:chord}
 d\epsilon_c=
 2\sum_{a<b}\sigma_{ca}\sigma_{cb}u_au_b
 +h\sum_a(\sigma_{ca}-\sigma_{ca}^2)u_a.
\end{equation}
Its term in degree $-j$ is free of rank
$2^k\binom\beta j$ over $R$. The total rank is $2^{|E(G)|}$.
Define the finite set
\[
 \mathcal A_G=
 \left\{s\in\{0,1\}^k:
       \sum_a\sigma_{ca}s_a\in\{0,1\}
       \text{ for every chord }c\right\}.
\]
Over $K=\mathbb C((h))$ there is a multiplicative
quasi-isomorphism
\begin{equation}\label{ei:esf:eq:graph-generic}
 C_{G,h}\otimes_R K\longrightarrow
 \prod_{s\in\mathcal A_G}
       \Lambda_K(\epsilon_1,\ldots,\epsilon_\beta),
 \qquad d=0.
\end{equation}
\end{theorem}

\begin{proof}
Tree differences are integral unimodular coordinates
on the translation quotient. Their monic quadratics
are a regular sequence. Quotient by them and send the
tree Koszul generators to zero, retaining the chord
generators. This is a multiplicative quasi-isomorphism,
as in Theorem~\ref{ei:esf:thm:flat-triangle}.

Expanding $\ell_c^2+h\ell_c$ and substituting
$u_a^2=-hu_a$ gives \eqref{ei:esf:eq:chord}.
In particular a negatively traversed tree edge contributes
$-2hu_a$, whose linear term cannot be discarded.
The squarefree monomials in the $u_a$, tensored with
exterior monomials, give the asserted ranks.

After inverting $h$, the tree algebra splits into the
points $u_a=-hs_a$, $s_a\in\{0,1\}$. On that factor the
chord differential is
$h^2L_c(s)(L_c(s)-1)$, where $L_c(s)=\sum_a\sigma_{ca}s_a$.
If one such scalar is nonzero it is invertible in $K$.
Multiplication by $\epsilon_c/(h^2L_c(s)(L_c(s)-1))$
contracts that factor. On the remaining factors every
chord differential is zero. Projection onto their product
is \eqref{ei:esf:eq:graph-generic}.
\end{proof}

Cohomology over $R$ distinguishes graphs with the same
generic description. Write $K_n$ for the complete graph
on $n$ vertices, $ij$ for the edge $\{i,j\}$, and
$C_{G,0}=C_{G,h}\otimes_R R/(h)$.
In the following table the dimensions are listed in degrees
$0,-1,\ldots$. Torsion refers to cohomology modules over $R$.

\begin{proposition}\label{ei:esf:prop:graph-examples}
For the graphs specified below, the central dimensions,
generic dimensions and torsion submodules are
\[
\begin{array}{c|c|c|c}
 G&H^*(C_{G,0})&H^*(C_{G,h}\otimes_RK)&R\text{-torsion}\\\hline
 K_3&(3,3)&(3,3)&0\\
 01,12,23,03&(6,6)&(4,4)&(R/(h))^2\text{ in }H^0\\
 K_4\setminus\{01\}&(5,10,5)&(5,10,5)&0\\
 K_4&(4,14,14,4)&(4,12,12,4)&(R/(h))^2\text{ in }H^{-1}.
\end{array}
\]
In the third row the colors of vertices $0,1,2,3$ are
$(1,1,2,3)$, and only intercolor pair factors are included.
The last row is an auxiliary four-color graph.
\end{proposition}

\begin{proof}
Use the following trees and chord path vectors, in the
orders shown:
\[
\begin{array}{c|c|c}
 G&T& (\sigma_c)_c\\\hline
 K_3&(01,12)&(1,1)\\
 01,12,23,03&(01,03,23)&(-1,1,-1)\\
 K_4\setminus\{01\}&(02,12,13)&(1,-1,1),\ (0,-1,1)\\
 K_4&(01,12,23)&(1,1,0),\ (1,1,1),\ (0,1,1).
\end{array}
\]
There are respectively $3,4,5,4$ admissible assignments.
Theorem~\ref{ei:esf:thm:graphs} gives all generic dimensions.

For the calculation over $R$, order squarefree monomials
by the bit mask of their tree coordinates, with $u_1$
the least significant bit. Order exterior subsets
lexicographically. Multiplication by $u_a$ is the tensor
product of identity matrices with
$\left(\begin{smallmatrix}0&0\\1&-h\end{smallmatrix}\right)$
in the $a$th position. Equation~\eqref{ei:esf:eq:chord}
therefore specifies every differential matrix. Removing
the $r$th exterior factor has sign $(-1)^{r-1}$.

Cancel a constant nonzero entry $w$ at row $r$, column
$c$, replacing the other entries by
\[
 D'_{ij}=D_{ij}-D_{ic}w^{-1}D_{rj}
 \quad(i\ne r,\ j\ne c).
\]
Apply the inverse column change to the preceding
differential and the inverse row change to the following
one. Their components in the cancelled summand vanish
because consecutive differentials compose to zero.
Thus each cancellation splits off the acyclic complex
$[R\xrightarrow{1}R]$ over $R$.

Choose the first available constant entry, reading
differentials from lowest degree upwards and each matrix
by rows. There are $1,2,6,14$ such cancellations in the
four cases. For the triangle and the third graph the
remaining differentials are zero, with respective term
ranks $(3,3)$ and $(5,10,5)$.

For the four-cycle the residual map has size $6$ by $6$.
Its only nonzero rows, after factoring out $h$, are
\[
 \begin{pmatrix}
 -2h&2h^2&4h&-4h^2&2h^2&-2h^3\\
 -2&2h&-2&2h&2h&-2h^2\\
 -2&2h&4&-4h&2h&-2h^2
 \end{pmatrix}.
\]
The first row is $h$ times the third. The other two
rows have a constant invertible minor in columns $1,3$.
Every remaining column is an $R$-linear combination
of these. Hence the two nonzero Smith factors of the
original residual map are $h,h$.

For $K_4$ the remaining term ranks are $(4,14,14,4)$.
The sole nonzero differential is from degree $-2$ to
degree $-1$. After factoring out $h$, deleting its zero
rows and columns gives
\[
 \begin{pmatrix}
 -2&2h&2&0&-2h&-2h&0\\
 2&-2h&0&-2&2h&0&2h\\
 0&0&-2&2&0&2h&-2h
 \end{pmatrix}.
\]

The rows sum to zero, and the first two rows have a
constant invertible minor in columns $1,3$. All other
columns lie in their span. Again the nonzero Smith
factors are $h,h$. The cokernel therefore has two
$R/(h)$ summands, now in degree $-1$.
The residual differentials vanish modulo $h$, giving
the central dimensions in the table. This also shows
that no higher power of $h$ occurs in the torsion.
\end{proof}

The absence of cohomological torsion for the triangle
is consistent with Proposition~\ref{ei:esf:prop:torsion-class}:
that proposition concerns a Hochschild class of its
multiplication. In the repeated-color row a full Hall
coefficient must also include the same-color factor.
The number of admissible graph branches by itself
defines neither a compact BPS multiplicity nor an Igusa
coefficient.


\section{The family over the spectral elliptic curve}
\label{ei:sec:global-family}

The formal calculation has a global source. On
$E_x\times E_y\times E_h$ put $u_1=x$, $u_2=y$, $u_3=x+y$,
and let
\[
 \mathcal S_i=u_i^*\Theta\otimes(u_i+h)^*\Theta,
 \qquad s_i=\vartheta(u_i)\vartheta(u_i+h).
\]
Define
\begin{equation}\label{ei:esf:eq:global}
 \mathscr K=\bigotimes_{i=1}^3
    [\mathcal S_i^{-1}\xrightarrow{s_i}\mathcal O],
 \qquad
 \mathscr C=R\pi_*\mathscr K,
 \quad\pi:E_x\times E_y\times E_h\to E_h.
\end{equation}
The displayed factors have degrees $-1,0$.
Write $E[n]=\ker([n]:E\to E)$ for the $n$-torsion subgroup.

\begin{theorem}\label{ei:esf:thm:global-fibers}
The complex $\mathscr C$ is represented by a two-term
complex of locally free sheaves of rank four on $E_h$.
Its derived fiber dimensions are
\[
 \begin{array}{c|cc}
  h&\dim H^0&\dim H^{-1}\\\hline
  0&3&3\\
  h\ne0,\ 2h\ne0&3&3\\
  h\ne0,\ 2h=0&4&4.
 \end{array}
\]
Near a nonzero point $h_0\in E[2]$, three summands have
zero differential. The fourth has local form
$[\mathcal O\xrightarrow{v}\mathcal O]$, with $v$ a
parameter at $h_0$ times a unit.
It contributes a skyscraper to $H^0(\mathscr C)$.
Its extra negative fiber class is a $\operatorname{Tor}_1$ class.
\end{theorem}

\begin{proof}
The divisors of $s_1,s_2$ cut out
\[
 Z=\{x=0\text{ or }-h\}\times_{E_h}
   \{y=0\text{ or }-h\},
\]
with their Cartier multiplicities. Each factor is finite
flat of degree two over $E_h$: the two sections have no
vertical component, and their sum has constant degree two.
Equivalently, near their meeting point its equation is a
monic quadratic in a relative local coordinate. Their
product is finite flat of degree four. The first two
sections are a regular sequence and their Koszul complex
resolves $\mathcal O_Z$.

Consequently \eqref{ei:esf:eq:global} reduces to the finite
pushforward of
$[\mathcal S_3^{-1}|_Z\xrightarrow{s_3}\mathcal O_Z]$.
An invertible sheaf on a finite flat degree-four scheme
has a locally free rank-four pushforward. Both terms are
therefore locally free, and specialization computes every
derived fiber.

For $h\ne0$, the four points of $Z_h$ are distinct.
The third section is zero at $(0,0),(0,-h),(-h,0)$.
At $(-h,-h)$ its value, in the corresponding line, is
\begin{equation}\label{ei:esf:eq:fourth-branch}
 \vartheta(-2h)\vartheta(-h).
\end{equation}
It is nonzero unless $2h=0$.

At a nonzero two-torsion
point the first factor has a simple zero, since $[2]$
is etale, and the second is nonzero. Local line frames
give the fourth complex in the statement. Its cokernel
is $\mathcal O/(v)$ and its kernel is zero. Derived
specialization creates the extra class in degree $-1$.

At $h=0$, formal logarithms and invertible theta frames
give \eqref{ei:esf:eq:triangle}, with dimensions $(3,3)$.
\end{proof}

For $h\ne0$ the fiber is formal: each surviving reduced
tree point carries one exterior factor with zero
differential, and a nonsurviving point is contractible.
At the three nonzero two-torsion points there are four
such exterior factors. No trivialization of the three
persistent line summands over the parameter curve is
needed for this local description.

On the original cover $h=2t$, the exceptional extra fibers
lie at $E[4]\setminus E[2]$, a set of twelve points.
The nonformal doubled fiber lies at the four points of
$E[2]$. A formal neighborhood of $t=0$ does not contain
the twelve exceptional points.

\subsection{A root line on the halving cover}

Choose a rigidification of the symmetric line $\Theta$
at the origin. Regard
$\mathcal S=u^*\Theta\otimes(u+h)^*\Theta$ as a line
on $E_u\times E_h$. For each fixed $h$ it has degree two
on $E_u$, so it has four degree-one square roots.
Their variation is a separate question.

\begin{theorem}\label{ei:esf:thm:root-cover}
There is no line $\mathcal R$ on $E_u\times E_h$ with
$\mathcal R^{\otimes2}\cong\mathcal S$.
There is also no continuous relative choice of a
degree-one root, even after tensoring $\mathcal S$ by
a line pulled back from $E_h$.
On the cover $h=2t$ there is an isomorphism
\begin{equation}\label{ei:esf:eq:parallelogram}
 u^*\Theta\otimes(u+2t)^*\Theta
 \cong (u+t)^*\Theta^{\otimes2}\otimes t^*\Theta^{\otimes2},
\end{equation}
and hence the root
\begin{equation}\label{ei:esf:eq:root}
 \mathcal R_t=(u+t)^*\Theta\otimes t^*\Theta.
\end{equation}
For every nonzero $\xi\in E[2]$, the deck transformation
$t\mapsto t+\xi$ changes its restriction to a $u$-fiber
by the nontrivial two-torsion line associated to $\xi$.
For $2t\ne0$, the section
$\vartheta(u)\vartheta(u+2t)$ is not the square of a
regular section of $\mathcal R_t$.
\end{theorem}

\begin{proof}
Restrict $\mathcal S$ to $\{u_0\}\times E_h$. Its first
factor is constant and its second has degree one.
The square of a line has even degree, which proves the
first assertion. Under $\operatorname{Pic}^1(E)\cong E$,
a root on a fixed fiber has a point $m$ with $2m=-h$.
A continuous choice over $E_h$ would induce
$2m_*=-\mathrm{id}$ on $\pi_1(E)=\mathbb Z^2$, which is
impossible. A line pulled back from the base does not
change these relative Picard classes.

For a symmetric rigidified line the cube theorem gives
the parallelogram identity
\[
 (f+g)^*\Theta\otimes(f-g)^*\Theta
 \cong(f^*\Theta)^{\otimes2}\otimes(g^*\Theta)^{\otimes2}
\]
for maps $f,g$ to $E$. The rigidification removes the
constant fiber factor. Apply the identity with $f=u+t$, $g=t$
to obtain \eqref{ei:esf:eq:parallelogram}.\footnote{
J.~S.~Milne, \emph{Abelian Varieties},
\href{https://www.jmilne.org/math/CourseNotes/AV.pdf}{Section~5},
the cube theorem and its symmetric-line consequences.
Stacks Project, \href{https://stacks.math.columbia.edu/tag/0BF4}{
Theorem of the cube, tag 0BF4}.}
This proves \eqref{ei:esf:eq:root} with its square isomorphism.

On a $u$-fiber the ratio under a deck shift is the translate
of $T_\xi^*\Theta\otimes\Theta^{-1}$, where $T_\xi$ is
translation by $\xi$. The degree-one polarization
$E\to\operatorname{Pic}^0(E)$ is an isomorphism, so this
line is nontrivial for $\xi\ne0$ and has order two.

Finally the two zeros $0,-2t$ of the displayed section
are distinct and simple when $2t\ne0$. A square of a
regular section has even vanishing order at every point.
\end{proof}

The isomorphism \eqref{ei:esf:eq:parallelogram} is a line
identity. The scalar equation
$\vartheta(u)\vartheta(u+2t)=
(\vartheta(u+t)\vartheta(t))^2$ is false:
its right-hand side vanishes identically at $t=0$.

For an ordered flag, tensor \eqref{ei:esf:eq:root} over its
selected interacting cross-pairs, using the same fixed
square isomorphism for every labeled pair. Every
parenthesization contains each pair exactly once.
Associativity and symmetry of tensor products of lines
therefore give all flag coherences, and their squares
are the corresponding coefficient-line identifications.
For the negative kernel lines use the inverse roots.
These constructions concern the specified arrow factors.
Same-color and dynamical factors require their own lines.

\section{The compact coefficient comparison}
\label{ei:sec:comparison}

The pair map, its simultaneous derived intersection, and its
parameter module retain different data. A multiplicative coefficient
comparison must specify which of these objects is its target.
For the graph intersection it must also preserve the higher products,
the trace character and derived specialization.
The independent quadratic family cannot replace the actual elliptic
displacement family by a change of base.

Let $S$ be a smooth projective K3 surface and let
$E_{\mathrm{geom}}$ be an elliptic curve.
A reduced critical Hall theory on $S\times E_{\mathrm{geom}}$
requires its own compact moduli, orientations and extension maps.
Extension flags retain their common translation torsor.
Their empty object lies over $BE_{\mathrm{geom}}$,
the classifying stack of these torsors.
The quotient of nonempty spectral configurations constructed above
does not identify this geometric base.

An operator representation has a different preservation requirement.
Multiplication by $t$ on the cochain complex $\mathbb C[t]$, with
zero differential, is a cochain endomorphism.
It is neither a unital algebra homomorphism nor a derivation.
Thus an action on a state complex is governed by its endomorphism
algebra. It need not act by automorphisms of an auxiliary
multiplication on that state space.

\begin{question}\label{ei:qn:compact}
Construct the reduced oriented critical coefficient functor for
the common-torsor extension geometry of $S\times E_{\mathrm{geom}}$.
Specify its extension pullbacks and supported pushforwards.
Give a map to a spectral base and determine its recipient:
the product coefficient cone, the simultaneous pair intersection,
or an independently defined state complex.

Can a multiplicative comparison retain the ordered coefficient lines,
their ambient images, the class $[m_3]$ and
$\delta_HG_2=hm_3$? Can it extend over the whole spectral elliptic
curve while preserving the exceptional fibers and the halving cover
with its action on root choices?
For a state representation, construct the corresponding supported
operators and identify its marked reduced state.
\end{question}

A regular multiplicative comparison cannot replace $G_2$ by $G_2/h$
without changing the coefficient ring.
The cyclic trace of the spectral model also requires a comparison
with compact reduced integration.
The latter has its own dual-number-valued state, as expressed by the
off-diagonal module coefficient in Section~\ref{nd:sec:comparison}.
The definition of reduced integration retains this distinction
already at the numerical level.\footnote{G.~Oberdieck and J.~Shen,
\emph{Reduced Donaldson--Thomas invariants and the ring of dual numbers},
\href{https://arxiv.org/abs/1612.03102}{Introduction}.
See the definition of the reduced integration map.}

Higher comparison terms must be constructed in the actual supported
Hom complex. The bar residual criterion at the end of
Chapter~\ref{nd:chap:diagonal} applies once that complex and its
cohomology projection are supplied.
The finite incidence contractions compute the operations of their
own coefficient algebras.

Finally, the root line on a fixed elliptic curve is not yet a
weight-five automorphic line. The Igusa comparison requires a
genus-two modular family and an identification of its geometric line
and section with the denominator line and its automorphic section.
The local theta factors and the chosen graph trace specify inputs
for this construction, rather than the required modular comparison.
